

\bigskip



{

This section describes how to tune Dyninst for optimum performance. During the course of a run, Dyninst will perform

several types of analysis on the binary, make safety assumptions about instrumentation that is inserted, and rewrite

the binary (perhaps several times). Given some guidance from the user, Dyninst can make assumptions about what work

it needs to do and can deliver significant performance improvements.}





\bigskip



{

There are two areas of Dyninst performance users typically care about. First, the time it takes Dyninst to parse and

instrument a program. This is typically the time it takes Dyninst to start and analyze a program, and the time it

takes to modify the program when putting in instrumentation. Second, many users care about the time instrumentation

takes in the modified mutatee. This time is highly dependent on both the amount and type of instrumentation put it,

but it is still possible to eliminate some of the Dyninst overhead around the instrumentation.}





\bigskip



{

The following subsections describe techniques for improving the performance of these two areas.}





\bigskip



\subsubsection{Optimizing Mutator Performance}



\bigskip



{

CPU time in the Dyninst mutator is usually consumed by either parsing or instrumenting binaries. When a new binary is

loaded, Dyninst will analyze the code looking for instrumentation points, global variables, and attempting to identify

functions in areas of code that may not have symbols. Upon user request, Dyninst will also parse debug information

from the binary, which includes local variable, line, and type information.





\bigskip



{

Since Dyninst 10.0.0, Dyninst supports parsing binaries in parallel, which significantly improve the analysis speed. We

typically have about 4X speedup when analyzing binaries with 8 threads. By default, Dyninst will use all the available

cores on your system. Please set environment variable OMP_NUM_THREADS to the number of desired threads.





\bigskip



{

Debugging information is lazily parsed separately from the rest of the binary parsing. Accessing line, type, or local

variable information will cause Dyninst to parse the debug information for all three of these.}





\bigskip



{

Another common source of mutator time is spent re-writing the mutatee to add instrumentation. When instrumentation is

inserted into a function, Dyninst may need to rewrite some or all of the function to fit the instrumentation in. If

multiple pieces of instrumentation are being inserted into a function, Dyninst may need to rewrite that function

multiple times.}





\bigskip



{

If the user knows that they will be inserting multiple pieces of instrumentation into one function, they can batch the

instrumentation into one bundle, so that the function will only be re-written once, using the

BPatch_process::beginInsertionSet and BPatch_­process::end­Inser­tion­Set functions (see section

\ref{bkm:Ref160339594}). Using these functions can result in a significant performance win when inserting

instrumentation in many locations.}





\bigskip



{

To use the insertion set functions, add a call to beginInsertionSet before inserting instrumentation. Dyninst will

start buffering up all instrumentation insertions. After the last piece of instrumentation is inserted, call

finalizeInsertionSet, and all instrumentation will be atomically inserted into the mutatee, with each function being

rewritten at most once.}





\bigskip



\subsubsection{Optimizing Mutatee Performance}



\bigskip



{

As instrumentation is inserted into a mutatee, it will start to run slower. The slowdown is heavily influenced by

three factors: the number of points being instrumented, the instrumentation itself, and the Dyninst overhead around

each piece of instrumentation. The Dyninst overhead comes from pieces of protection code (described in more detail

below) that do things such as saving/restoring registers around instrumentation, checking for instrumentation

recursion, and performing thread safety checks.





\bigskip



{

The factor by which Dyninst overhead influences mutatee run-time depends on the type of instrumentation being inserted.

When inserting instrumentation that runs a memory cache simulator, the Dyninst overhead may be negligible. On the

other-hand, when inserting instrumentation that increments a counter, the Dyninst overhead will dominate the time spent

in instrumentation. Remember, optimizing the instrumentation being inserted may sometimes be more important than

optimizing the Dyninst overhead. Many users have had success writing tools that make use of

Dyninst's ability to dynamically remove instrumentation as a performance improvement.}





\bigskip



{

The instrumentation overhead results from safety and correctness checks inserted by Dyninst around instrumentation.

Dyninst will automatically attempt to remove as much of this overhead as possible, however it sometimes must make a

conservative decision to leave the overhead in. Given additional, user-provided information Dyninst can make better

choices about what safety checks to leave in. An unoptimized post-Dyninst 5.0 instrumentation snippet looks like the

following:}





\bigskip



\begin{flushleft}

\tablefirsthead{}

\tablehead{}

\tabletail{}

\tablelasttail{}

\begin{supertabular}{|m{3.24696in}|m{3.25186in}|}

\hline

{ Save General Purpose Registers} &

{ In order to ensure that instrumentation doesn't corrupt the program,

Dyninst saves all live general purpose registers.}\\\hline

{ Save Floating Point Registers} &

{ Dyninst may decide to separately save any floating point registers that may be corrupted by

instrumentation.}\\\hline

{ Generate A Stack Frame} &

{ Dyninst builds a stack frame for instrumentation to run under. This provides the illusion to

instrumentation that it is running as its own function.}\\\hline

{ Calculate Thread Index} &

{ Calculate an index value that identifies the current thread. This is primarily used as input

to the Trampoline Guard.}\\\hline

{ Test and Set Trampoline Guard} &

{ Test to see if we are already recursively executing under instrumentation, and skip the user

instrumentation if we are.}\\\hline

{ Execute User Instrumentation} &

{ Execute any BPatch_snippet code.}\\\hline

{ Unset Trampoline Guard} &

{ Marks the this thread as no longer being in instrumentation}\\\hline

{ Clean Stack Frame} &

{ Clean the stack frame that was generated for instrumentation.}\\\hline

{ Restore Floating Point Registers} &

{ Restore the floating point registers to their original state.}\\\hline

{ Restore General Purpose Registers} &

{ Restore the general purpose registers to their original state.}\\\hline

\end{supertabular}

\end{flushleft}

{

Dyninst will attempt to eliminate as much of its overhead as is possible. The Dyninst user can assist Dyninst by doing

the following:}





\bigskip



\begin{itemize}

\item {

Write BPatch_snippet code that avoids making function calls. Dyninst will attempt to perform analysis on the user

written instrumentation to determine which general purpose and floating point registers can be saved. It is difficult

to analyze function calls that may be nested arbitrarily deep. Dyninst will not analyze any deeper than two levels of

function calls before assuming that the instrumentation clobbers all registers and it needs to save everything.}

\end{itemize}

{

In addition, not making function calls from instrumentation allows Dyninst to eliminate its tramp guard and thread index

calculation. Instrumentation that does not make a function call cannot recursively execute more instrumentation.

}



\begin{itemize}

\item {

Call BPatch::setTrampRecursive(true) if instrumentation cannot execute recursively. If instrumentation must make a

function call, but will not execute recursively, then enable trampoline recursion. This will cause Dyninst to stop

generating a trampoline guard and thread index calculation on all future pieces of instrumentation. An example of

instrumentation recursion would be instrumenting a call to write with instrumentation that calls

printf$\text{\textgreek{—}}$write will start calling printf printf will re-call write.}

\item {

Call BPatch::setSaveFPR(false) if instrumentation will not clobber floating point registers. This will cause

Dyninst to stop saving floating point registers, which can be a significant win on some platforms.}

\item {

Use simple BPatch_snippet objects when possible. Dyninst will attempt to recognize, peep-hole optimize, and simplify

frequently used code snippets when it finds them. For example, on x86 based platforms Dyninst will recognize snippets

that do operations like $\text{\textgreek{‘}}$var = constant' or

$\text{\textgreek{‘}}$var++' and turn these into optimized assembly instructions that take

advantage of CISC machine instructions.}

\item {

Call BPatch::setInstrStackFrames(false) before inserting instrumentation that does not need to set up stack frames.

Dyninst allows you to force stack frames to be generated for all instrumentation. This is useful for some

applications (e.g., debugging your instrumentation code) but allowing Dyninst to omit stack frames wherever possible

will improve performance. This flag is false by default; it should be enabled for as little instrumentation as

possible in order to maximize the benefit from optimizing away stack frames.}

\item {

Avoid conditional instrumentation wherever possible. Conditional logic in your instrumentation makes it more difficult

to avoid saving the state of the flags.}

\item {

Avoid unnecessary instrumentation. Dyninst provides you with all kinds of information that you can use to select only

the points of actual interest for instrumentation. Use this information to instrument as selectively as possible.

The best way to optimize your instrumentation, ultimately, is to know a priori that it was unnecessary and not insert

it.}

\end{itemize}
